\section{Related Work}
\label{sec:related}

\paragraph{Spiking Transformers.}
Spiking Transformers combine the sequence modeling capacity of Transformers
with the event-driven computation of SNNs. Spikformer introduced spike-form
self-attention for directly trained spiking Transformers~\cite{zhou2023spikformer}.
Spike-driven Transformer removes floating-point multiplication from spike-form
attention~\cite{yao2023spikedriven}, and SpikingResformer combines spiking
attention with a multistage convolutional hierarchy~\cite{shi2024spikingresformer}.
QKFormer further develops this line with hierarchical spike-form Q-K attention
based on binary token/channel vectors~\cite{zhou2024qkformer}. Our work is not a
replacement for these architectures; it asks whether QKFormer's existing
binary states form useful, auditable lookup addresses.

\paragraph{Transformer-to-SNN conversion.}
ANN-to-SNN conversion has been successful for convolutional networks, but
Transformer attention introduces additional difficulty because token
interactions depend on softmax normalization and continuous query--key scores.
Recent conversion work such as SpikedAttention directly addresses attention
conversion for Transformers~\cite{hwang2024spikedattention}. In contrast, this
paper begins from a spiking Transformer whose attention already exposes binary
Q/K structure, and asks whether that structure can index auditable lookup
prototypes and controlled downstream probes.

\paragraph{Lookup and compositional memory.}
Learned lookup has been used to replace arithmetic with centroid-indexed table
access in MADDNESS and LUT-NN~\cite{blalock2021maddness,tang2023lutnn}, while
Product-Key Memory factorizes a large key space into compositional
subkeys~\cite{lample2019large}. Product quantization decomposes a vector space
into a Cartesian product of lower-dimensional codebooks, while additive
quantization reconstructs vectors as sums of codewords from multiple
codebooks~\cite{jegou2011product,babenko2014additive}. LookupFFN and MemoryFormer study table-based
replacements for Transformer feed-forward
computation~\cite{zeng2023lookupffn,wang2024memoryformer}. Our setting differs in both target and
claim: we audit binary states already present inside spike-form Q-K attention
and reconstruct local module responses rather than replacing dense ANN layers.
The subspace QK-LUT is closest in spirit to compositional memory and
multi-codebook quantization, but it partitions an existing discrete address by
QKFormer semantics and stores conditional response residuals rather than
quantizing input vectors. It is also evaluated against same-entry shuffled
controls. Because lookup can still lose input-dependent context, we treat
downstream replacement as a separate gate rather than infer it from
reconstruction alone.

\begin{table*}[t]
\centering
\scriptsize
\setlength{\tabcolsep}{4pt}
\renewcommand{\arraystretch}{1.15}
\resizebox{\textwidth}{!}{%
\begin{tabular}{p{0.18\textwidth}p{0.21\textwidth}p{0.19\textwidth}p{0.19\textwidth}p{0.25\textwidth}}
\toprule
Research direction & Representative work & Address or state source &
Primary objective & Difference from QK-LUTFormer \\
\midrule
Spiking attention architectures &
Spikformer, Spike-driven Transformer, SpikingResformer,
QKFormer~\cite{zhou2023spikformer,yao2023spikedriven,shi2024spikingresformer,zhou2024qkformer} &
Spikes and architecture-specific Q/K or SSA states &
Design accurate, spike-driven attention backbones &
They provide the binary states and interaction structures that we audit; they
do not test Q/K address semantics with matched coarse and shuffled controls. \\

Transformer-to-SNN conversion &
SpikedAttention~\cite{hwang2024spikedattention} &
Converted ANN activations and attention states &
Preserve ANN accuracy after temporal conversion &
We begin from natively spiking Q/K states and study lookup addressability and
replacement fidelity rather than ANN-to-SNN conversion. \\

Arithmetic LUT inference &
MADDNESS, LUT-NN~\cite{blalock2021maddness,tang2023lutnn} &
Learned centroids or quantized activation codes &
Replace dense arithmetic with table access &
Our indices are deterministic semantic Q/K spike fields, and the central test
is correct address--response alignment rather than arithmetic replacement
alone. \\

Compositional codebooks and memory &
Product-Key Memory, product and additive
quantization~\cite{lample2019large,jegou2011product,babenko2014additive} &
Learned subkeys or continuous-vector codebooks &
Scale retrieval capacity or compress vector reconstruction &
Our subtables follow QKFormer token/channel, Q/gate, and K semantics, include
support-aware fallback, and are evaluated against matched shuffled subtables. \\

Transformer lookup modules &
LookupFFN, MemoryFormer~\cite{zeng2023lookupffn,wang2024memoryformer} &
Token, hidden-state, or feed-forward lookup keys &
Replace or augment Transformer feed-forward computation &
We target spiking-attention responses and explicitly test temporal/channel
calibration, current replacement, and native affine/BN/LIF operator fidelity. \\

\textbf{QK-LUTFormer / TCSLU (ours)} &
\textbf{Semantic-address audit and calibrated replacement} &
\textbf{Binary Q/K, token/channel, and gate states} &
\textbf{Compact auditable lookup with bounded replacement fidelity} &
\textbf{Combines negative-controlled address semantics, explicit misses,
subspace compression, dynamic calibration, and qualified operator-class
replacement in one evidence chain.} \\
\bottomrule
\end{tabular}
}
\caption{Positioning against adjacent research areas. The table compares the
source of lookup indices and the scientific question being tested, rather than
ranking unlike datasets or hardware platforms. QK-LUTFormer does not claim to
invent LUT inference or spiking attention; it connects semantic Q/K addressing,
matched negative controls, compact subtable composition, and calibrated
replacement fidelity.}
\label{tab:related-work-positioning}
\end{table*}
